% TOP_PROBLEM: {"id": 3379, "title": "Sets with distinct subset sums", "category_id": 1, "category": {"id": 1, "name": "number_theory", "display_name": "Number Theory", "description": "Properties of integers, prime numbers, Diophantine equations.", "slug": "number-theory", "order_index": 1, "created_at": "2026-07-31T15:26:25.670Z"}, "status": "open"}
% FIRST_POSTED: null
% ENTRY_KIND: statement_only
% Imported from: batch04.tex; UnsolvedMath ID 3379
\documentclass[11pt]{article}
\usepackage[margin=25mm]{geometry}
\usepackage{fontspec}
\setmainfont{Latin Modern Roman}
\usepackage{amsmath,amssymb,mathtools}
\usepackage{unicode-math}
\setmathfont{Latin Modern Math}
\usepackage{xurl,hyperref}
\hypersetup{colorlinks=true,urlcolor=blue}
\setcounter{secnumdepth}{0}
\setlength{\parindent}{0pt}
\setlength{\parskip}{5pt}
\setlength{\emergencystretch}{3em}
\newcommand{\sref}[2]{\hyperlink{source:#1:#2}{[S#2]}}
\newcommand{\cataloguescope}[1]{\paragraph{Recorded target.}#1}
\begin{document}
% UnsolvedMath ID 3379; source catalogue code OPG-414
\setcounter{section}{3378}
\section{Sets with distinct subset sums}\label{3379}
{\small\sffamily UnsolvedMath ID 3379 · Number Theory}
\subsection{Definitions and mathematical statement}

\paragraph{Shared notation.}$\mathbb N=\{1,2,\ldots\}$, and $\mathbb Z,\mathbb Q,\mathbb R,\mathbb C$ denote the integers, rationals, reals and complex numbers. Any other conventions are specified locally.


A finite set $S\subseteq\mathbb N$ has distinct subset sums if the function $T\mapsto\sum_{s\in T}s$ on $\mathcal P(S)$ is injective; the empty sum is zero. For $m\ge1$, put
\[
 f(m)=\min\{\max S:S\subseteq\mathbb N,\ |S|=m,\ S\text{ has distinct subset sums}\}.
\]
\paragraph{Statement recorded in the supplied source.}
Is there an absolute constant $C\in\mathbb R$ such that, for every $N\ge1$ and every sum-distinct $S\subseteq\{1,\ldots,N\}$,
\[
 |S|\le\log_2 N+C?
\]
Equivalently, is there a constant $c>0$ with $f(m)\ge c2^m$ for all $m\ge1$? More generally, determine the asymptotic behavior of $f(m)$. \sref{3379}{2}
\subsection{Short English statement}
Bound the size of a sum-distinct set by the binary logarithm of its largest permitted element plus an absolute constant, and determine the corresponding minimum-largest-element function.
\subsection{Sources}
\begin{description}
\item[\hypertarget{source:3379:1}{S1}] ULAM.AI and the UnsolvedMath Contributors, UnsolvedMath: A Curated Collection of Open Mathematics Problems, record 3379 (OPG-414). CC BY 4.0. \url{https://huggingface.co/datasets/ulamai/UnsolvedMath}.
\item[\hypertarget{source:3379:2}{S2}] Open Problem Garden, Sets with distinct subset sums, OPG-414. Original problem statement, full mathematical discussion, bibliography and comments. \url{https://www.openproblemgarden.org/op/sets_with_distinct_subset_sums}.
\end{description}
\label{end:3379}
\end{document}
